\subsection{DERIVED SERIES, SOLVABLE GROUPS}

DEFINITION 7. Let G be a group. The derived series of G is the series \((\mathbf{D}^n (\mathbf{G}))_{n\in \mathbb{N}}\) defined inductively by:

\[
\mathbf {D} ^ {0} (\mathbf {G}) = \mathbf {G}; \quad \mathbf {D} ^ {n + 1} (\mathbf {G}) = \mathbf {D} (\mathbf {D} ^ {n} (\mathbf {G})) \quad \text { for } n \in \mathbf {N}.
\]

Then \(\mathbf{D}^0 (\mathbf{G}) = \mathbf{C}^1 (\mathbf{G}) = \mathbf{G},\mathbf{D}^1 (\mathbf{G}) = \mathbf{C}^2 (\mathbf{G}) = \mathbf{D}(\mathbf{G}) = (\mathbf{G},\mathbf{G}).\) For all \(n \in \mathbf{N}, \mathrm{D}^n(\mathrm{G}) \subset \mathrm{C}^{2^n}(\mathrm{G})\) , as is seen by induction on \(n\) using formula (7) of no. 3.

Let \(f: \mathbf{G} \to \mathbf{G}'\) be a group homomorphism. It is seen, by induction on \(n\) , that \(f(\mathbf{D}^n(\mathbf{G})) \subset \mathbf{D}^n(\mathbf{G}')\) and that, if \(f\) is surjective, \(f(\mathbf{D}^n(\mathbf{G})) = \mathbf{D}^n(\mathbf{G}')\) . In particular, for all \(n \in \mathbf{N}\) , \(\mathbf{D}^n(\mathbf{G})\) is a characteristic (and therefore normal) subgroup of \(\mathbf{G}\) . For all \(n \in \mathbf{N}\) , the group \(\mathbf{D}^n(\mathbf{G}) / \mathbf{D}^{n+1}(\mathbf{G})\) is a commutative normal (but not in general central) subgroup of \(\mathbf{G} / \mathbf{D}^{n+1}(\mathbf{G})\) .

Let \((\mathbf{G}_0, \mathbf{G}_1, \ldots)\) be a decreasing sequence of subgroups of \(\mathbf{G}\) such that: (1) \(\mathbf{G}_0 = \mathbf{G}\) ; (2) for all \(i\) , \(\mathbf{G}_{i+1}\) is normal in \(\mathbf{G}_i\) and \(\mathbf{G}_i / \mathbf{G}_{i+1}\) is commutative. Then \(\mathbf{D}^i(\mathbf{G}) \subset \mathbf{G}_i\) for all \(i\) , as is seen by induction on \(i\) .

DEFINITION 8. A group G is called solvable if there exists an integer n such that \(\mathbf{D}^{n}(\mathbf{G})=\{e\}\) . If G is a solvable group, the least integer n such that \(\mathbf{D}^{n}(\mathbf{G})=\{e\}\) is called the solvability class of G.

A solvable group of solvability class n is called a solvable group of class n.~A group is sometimes said to be of finite solvability class if it is solvable.

Examples. (1) A group is solvable of class 0 (resp. \(\leqslant 1\) ) if and only if it is reduced to \(\{e\}\) (resp. is commutative).

\begin{enumerate}
\def\labelenumi{(\arabic{enumi})}
\setcounter{enumi}{1}
\item
  Every nilpotent group of class \(\leqslant 2^{n} - 1\) is solvable of class \(\leqslant n\) ; this follows from the relation \(\mathbf{D}^n (\mathbf{G})\subset \mathbf{C}^{2^n}(\mathbf{G})\) proved above.
\item
  Let G be a solvable group of class \(\leqslant n\) . Every subgroup (resp. quotient group) of G is solvable of class \(\leqslant n\) (proof analogous to that of no. 3, Example 3).
\item
  If \(\mathbf{G}\) is a solvable group of class \(p\) and \(\mathbf{F}\) is a solvable group of class \(q\) , every extension \(\mathbf{E}\) of \(\mathbf{G}\) by \(\mathbf{F}\) is a solvable group of class \(\leqslant p + q\) . For, let \(\pi: \mathbf{E} \to \mathbf{G}\) be the projection; then \(\pi(\mathrm{D}^p(\mathbf{E})) \subset \mathrm{D}^p(\mathbf{G}) = \{e\}\) and therefore \(\mathrm{D}^p(\mathbf{E}) \subset \mathbf{F}\) ; it follows that \(\mathrm{D}^{p+q}(\mathbf{E}) = \mathrm{D}^q(\mathrm{D}^p(\mathbf{E})) \subset \mathrm{D}^q(\mathrm{F}) = \{e\}\) .
\item
  The symmetric group \(\mathfrak{S}_n\) is solvable if and only if \(n < 5\) (cf.~§ 5, Exercises 10 and 16).
\item
  *If A is a commutative ring, the upper triangular group T(n, A) is solvable but not in general nilpotent.*
\end{enumerate}

PROPOSITION 10. Let G be a group and n an integer. The following conditions are equivalent:

\begin{enumerate}
\def\labelenumi{(\roman{enumi})}
\item
  G is solvable of class \(\leqslant n\) .
\item
  There exists a series of normal subgroups of \(\mathbf{G}\)
\end{enumerate}

\[
\mathbf {G} = \mathbf {G} ^ {0} \supset \mathbf {G} ^ {1} \supset \dots \supset \mathbf {G} ^ {n} = \{e \}
\]

such that the groups \(\mathbf{G}^k /\mathbf{G}^{k + 1}\) are commutative.

\begin{enumerate}
\def\labelenumi{(\roman{enumi})}
\setcounter{enumi}{2}
\tightlist
\item
  There exists a series of subgroups of \(\mathbf{G}\)
\end{enumerate}

\[
\mathbf {G} = \mathbf {G} ^ {0} \supset \mathbf {G} ^ {1} \supset \dots \supset \mathbf {G} ^ {n} = e
\]

such that, for all \(k\) , \(\mathbf{G}^{k+1}\) is a normal subgroup of \(\mathbf{G}^k\) and \(\mathbf{G}^k/\mathbf{G}^{k+1}\) is commutative.

\begin{enumerate}
\def\labelenumi{(\roman{enumi})}
\setcounter{enumi}{3}
\tightlist
\item
  There exists a normal commutative subgroup \(\mathbf{A}\) of \(\mathbf{G}\) such that \(\mathbf{G} / \mathbf{A}\) is solvable of class \(\leqslant n - 1\) .
\end{enumerate}

For (i) \(\Rightarrow\) (ii) it suffices to take \(\mathbf{G}^k\) equal to \(\mathbf{D}^k (\mathbf{G})\) . (ii) \(\Rightarrow\) (iii) trivially. (iii) \(\Rightarrow\) (i) for \(\mathbf{D}^k (\mathbf{G})\) is necessarily contained in \(\mathbf{G}^k\) . The equivalence of (ii) and (iv) is immediate by induction on \(n\) .

More briefly: a group is solvable of class \(\leqslant n\) if it can be obtained by successive extensions of n commutative groups.

COROLLARY. Let G be a finite group and

\[
\mathbf {G} = \mathbf {G} ^ {0} \supset \mathbf {G} ^ {1} \supset \dots \supset \mathbf {G} ^ {n} = \{e \}
\]

a Jordan-Hölder series of G For G to be solvable, it is necessary and sufficient that the quotients \(G^{k}/G^{k+1}\) be cyclic of prime order.

If the quotients of a composition series of G are cyclic and hence commutative, G is solvable by Proposition 10. Conversely, if G is solvable, the group \(G^{k}/G^{k+1}\) is, for all k, solvable and simple (§ 4, no. 7, Proposition 9). Now, every solvable simple group H is cyclic of prime order. For D(H) is a normal subgroup of H; D(H) = H is impossible for in that case \(D^{k}(H) = H\) for all k; then D(H) = \{e\} and H is commutative. The corollary then follows from § 4, no. 10, Corollary to Proposition 20.

